\scanpage{13}
\begin{center}-- 54 --\end{center}
\hfill n° 371

\[
(17)\qquad \varepsilon(\alpha,\beta)=\prod_{i>j}\varepsilon_{ij}(\alpha_i,\beta_j).
\]

En effet, le second membre de (16) appartient évidemment à $E_{\alpha+\beta}$ et l'application $(x_1,\ldots,x_n,y_1,\ldots,y_n)\longrightarrow\varepsilon(\alpha,\beta)\bigotimes_{i\in I}(x_i y_i)$ est $A$-multilinéaire dans le produit des $E_{\alpha_i}$ et des $E_{\beta_j}$ $(1\leq i\leq n,\ 1\leq j\leq n)$. Il s'agit ensuite de prouver que la multiplication ainsi définie sur $E$ est \uline{associative}; or, si $\gamma=(\gamma_i)$ est un troisième élément de $\Delta$, et $z_i\in E_{i,\gamma_i}$ pour $1\leq i\leq n$, on a
\[
\begin{aligned}
((\bigotimes_i x_i)(\bigotimes_i y_i))(\bigotimes_i z_i)&=\varepsilon(\alpha+\beta,\gamma)\varepsilon(\alpha,\beta)\bigotimes_i(x_i y_i z_i)\\
(\bigotimes_i x_i)((\bigotimes_i y_i)(\bigotimes_i z_i))&=\varepsilon(\alpha,\beta+\gamma)\varepsilon(\beta,\gamma)\bigotimes_i(x_i y_i z_i)
\end{aligned}
\]
et tout revient donc à vérifier l'identité
\[
\varepsilon(\alpha+\beta,\gamma)\varepsilon(\alpha,\beta)=\varepsilon(\alpha,\beta+\gamma)\varepsilon(\beta,\gamma)
\]
qui résulte aussitôt des relations
\[
\begin{aligned}
\varepsilon(\alpha+\beta,\gamma)&=\varepsilon(\alpha,\gamma)\varepsilon(\beta,\gamma)\\
\varepsilon(\alpha,\beta+\gamma)&=\varepsilon(\alpha,\beta)\varepsilon(\alpha,\gamma)
\end{aligned}
\]
elles-mêmes conséquences immédiates de la définition (17).

Si $e_i$ désigne l'élément unité de $E_i$, on sait que $e_i$ est homogène de degré 0, donc $e=\bigotimes_{i\in I}e_i$ est homogène de degré 0, et il résulte de (16) et des relations $\varepsilon_{ij}(\alpha_i,0)=\varepsilon_{ij}(0,\beta_j)=1$ que $e$ est élément unité de $E$, ce qui achève de définir sur $E$ la structure de $A$-algèbre graduée cherchée. On prendra ensuite $h_i(x_i)=x_i\otimes\bigotimes_{j\ne i}e_j$; il résulte de (16) que $h_i$ est un homomorphisme d'algèbres et vérifie les conditions 1° et 2° de l'énoncé, en vertu des relations $\varepsilon_{ij}(0,\beta_j)=\varepsilon_{ij}(\alpha_i,0)=1$. Ces mêmes relations prouvent que l'on a
\[
(18)\qquad \bigotimes_{i\in I}=\prod_{i\in I}h_i(x_i)
\]
où le second membre est le produit de la \uline{séquence} $(h_i(x_i))_{i\in I}$ pour la structure
\mspage{12}
